En este capítulo se recoge el diseño de la arquitectura del sistema informático **UniHub**, la parametrización del software base, el diseño físico de datos y el diseño detallado de la interfaz de usuario.

\section{Arquitectura del Sistema}
Se adopta una arquitectura desacoplada basada en microservicios contenerizados y un patrón en capas (Layers) orientada a la web.

\subsection{Diseño de alto nivel}
Siguiendo el modelo C4 (Contenedores y Componentes), el sistema se divide en cuatro capas principales:
\begin{enumerate}
    \item \textbf{Capa de Presentación (Frontend):} Servidor Web Nginx entregando la SPA compilada en React con soporte para la ruta aislada \texttt{/admin} (Fase 3).
    \item \textbf{Capa de Servicio (API REST):} Framework FastAPI en Python que expone controladores RESTful, métricas cgroup de contenedores, sincronización reactiva \texttt{/api/v1/admin/sync-etl} y endpoints CRUD (Fase 2).
    \item \textbf{Capa de Datos (Persistencia):} Motor de Base de Datos Relacional PostgreSQL 15 con el esquema DDL y almacenamiento de archivos JSON en disco (Fase 2 y 1).
    \item \textbf{Capa de Ingesta (Crawler):} Módulo independiente en Python estructurado en 2 procesos paralelos y 3 partes (RUCT/BOE, webs oficiales y consolidación de precios ECTS SIIU) con servicio Cron para ejecuciones desatendidas mensuales (Fase 1).
\end{enumerate}

\subsection{Diseño detallado}
La API REST implementa controladores modulares mediante \texttt{APIRouter} en FastAPI:
\begin{itemize}
    \item \texttt{routes/universidades.py}: Endpoints con ordenación prioritaria (Públicas primero) y operaciones CRUD.
    \item \texttt{routes/titulaciones.py}: Endpoints para Grados, Másteres y Doctorados y lectura de planes de estudio BOE.
    \item \texttt{metrics/container\_metrics.py}: Muestreo de métricas físicas de salud (RAM RSS MB, CPU \%) y huella de carbono ($gCO_2$) por contenedor cgroup.
\end{itemize}

\section{Parametrización del software base}
Se configuró el servidor proxy inverso Nginx (\texttt{nginx.conf}) para redirigir las peticiones \texttt{/api/v1/} hacia la API REST en \texttt{http://api:8000/api/v1/} y servir la SPA React en la raíz \texttt{/}, soportando react-router para la URL \texttt{/admin}. En PostgreSQL, se parametrizó la creación del usuario restringido \texttt{unihub\_api\_user} con permisos \texttt{GRANT SELECT} exclusivo.

\section{Diseño Físico de Datos}
El diseño relacional en PostgreSQL (\texttt{schema.sql}) se compone de las siguientes tablas principales con claves primarias y foráneas con borrado en cascada, optimizadas mediante la extensión de trigramas \texttt{pg\_trgm}:
\begin{itemize}
    \item \texttt{universidades}: Clave primaria \texttt{codigo VARCHAR(10)}. Índices en \texttt{tipo}, \texttt{comunidad\_autonoma} e índice GIN trigrama \texttt{idx\_univ\_nombre\_trgm} sobre \texttt{nombre}.
    \item \texttt{titulaciones}: Clave primaria \texttt{codigo\_estudio VARCHAR(20)}, clave foránea \texttt{universidad\_codigo REFERENCES universidades(codigo)}. Índice GIN trigrama \texttt{idx\_tit\_titulo\_trgm} sobre \texttt{titulo}.
    \item \texttt{planes\_estudio}: Clave foránea \texttt{codigo\_estudio REFERENCES titulaciones(codigo\_estudio)}.
    \item \texttt{resumen\_creditos}: Clave foránea a \texttt{planes\_estudio(id)}.
    \item \texttt{elementos\_curriculares}: Clave foránea a \texttt{planes\_estudio(id)}. Tipos de datos en texto extenso (\texttt{TEXT}) para evitar truncamientos.
    \item \texttt{errores\_crawler} y \texttt{estadisticas\_rendimiento}: Tablas relacionales de registro de auditoría con desduplicación por timestamp.
\end{itemize}

\section{Diseño de la Interfaz de Usuario} 
El diseño visual se ha construido respetando la paleta de colores de la **Universidad de Cádiz (UCA)**:
\begin{itemize}
    \item \textbf{Azul Noche UCA (Principal):} \texttt{\#002B49} y \texttt{\#003366}.
    \item \textbf{Azul Mar de Cádiz (Acentos):} \texttt{\#0084C8} y \texttt{\#00A8E8}.
    \item \textbf{Dorado Sol de Cádiz (Highlights/Badges):} \texttt{\#F3A712} y \texttt{\#FFC72C}.
    \item \textbf{Rosa Magenta (Badge Doctorado):} \texttt{\#EC4899} (20\% opacidad).
    \item \textbf{Panel de Administración:} Vista multisección con monitor de métricas en tiempo real, tabla de diagnóstico del Checkpoint de la Fase 1, visor de errores JSON y pestaña interactiva con documentación Swagger UI (\texttt{/docs}) y ReDoc (\texttt{/redoc}).
    \item \textbf{Estilos Visuales:} Formularios con efecto desenfocado *glassmorphism* (\texttt{backdrop-filter: blur(12px)}), componentes de paginación modular, distancia integrada en tarjeta y modo oscuro dinámico.
\end{itemize}
